\section{附录}
\label{sec:appendix}

\subsection*{A. 21 个 MCP 工具完整清单}

按能力簇分组，与第 \ref{sec:product:mcp} 节的核心叙事对应。
所有工具通过 MCP 协议暴露给 agent，返回结构化 JSON 而非自然语言。

\begin{center}
\small
\begin{tabularx}{\linewidth}{@{}>{\tabRR}p{2.4cm}>{\tabRR}p{3.6cm}>{\tabRR}X@{}}
\toprule
\rowcolor{tableHead}
\heiti 能力簇 & \heiti 工具名 & \heiti 作用概要 \\
\midrule
\keyword{语义分层} & \code{layers\_split} & 根据 agent 提供的 orchestrated plan 把图像分解为多个带透明度的语义图层 \\
 & \code{layers\_list} & 列出已分解会话的 layers 元信息 \\
\midrule
\keyword{区域编辑} & \code{inpaint\_artwork} & 对指定 mask 区域重绘 \\
 & \code{layers\_edit} & 对指定 layer 原地编辑 \\
 & \code{layers\_redraw} & 对指定 layer 用新 prompt 重绘 \\
 & \code{layers\_transform} & 对 layer 做几何/色彩变换 \\
 & \code{layers\_composite} & 多 layer 合成 \\
 & \code{layers\_export} & 导出合成结果（PNG / 分层 PSD） \\
 & \code{layers\_evaluate} & 对合成结果做质量评估 \\
\midrule
\keyword{文化评估} & \code{evaluate\_artwork} & 基于 L1–L5 \& 指定传统做多维评价 \\
 & \code{list\_traditions} & 枚举已定义传统（当前 13 种） \\
 & \code{get\_tradition\_guide} & 获取指定传统的完整 YAML 规则 \\
 & \code{search\_traditions} & 按关键词搜索传统定义 \\
\midrule
\keyword{结构化生成} & \code{create\_artwork} & 基于 intent + tradition 生成新图像 \\
 & \code{generate\_image} & 通用文生图 \\
\midrule
\keyword{视觉 brief} & \code{brief\_parse} & 解析视觉意图，产出 proposal 所需字段 \\
 & \code{generate\_concepts} & 生成概念草图 / 选项 \\
\midrule
\keyword{运行管理} & \code{view\_image} & agent 把图像回传进 Claude Code 视野 \\
 & \code{unload\_models} & 释放 GPU/MPS 显存 \\
 & \code{archive\_session} & 归档会话工件 \\
 & \code{sync\_data} & 同步外部数据源（可选） \\
\bottomrule
\end{tabularx}
\end{center}

\subsection*{B. 13 种文化/设计传统评价体系示例}

每个传统的 YAML 都包含 L1–L5 判别规则、典型元素库、常见失败模式描述。
以下给出一个\keyword{中国工笔}的简化示例，
完整定义在 \code{src/vulca/cultural/data/traditions/chinese\_gongbi.yaml}。

\begin{quote}
\small\ttfamily
L1 形制：\\
\quad 画面以线造型，边界明确；\\
\quad 构图多采取"折枝式"或"边角式"，避免满幅。\\
L2 技法：\\
\quad 白描 / 重彩 / 淡彩 三种基础技法；\\
\quad 设色讲究层次叠染，每层需等待干透。\\
L3 语汇：\\
\quad 题材：花鸟、仕女、草虫、折枝；\\
\quad 符号：屏风、云水纹、印章。\\
L4 风格：\\
\quad 宋院体、明代"吴门"、近现代"于非闇"等流派差异。\\
L5 气韵：\\
\quad 讲究"形神兼备"；\\
\quad 评价时需引用 L1–L4 的具体证据，不做主观优劣判断。
\end{quote}

完整覆盖范围：
\code{african\_traditional}、\code{brand\_design}、\code{chinese\_gongbi}、\code{chinese\_xieyi}、
\code{contemporary\_art}、\code{islamic\_geometric}、\code{japanese\_traditional}、\code{photography}、
\code{south\_asian}、\code{ui\_ux\_design}、\code{watercolor}、\code{western\_academic}、\code{default}。

\subsection*{C. 开源与社区数据（持续更新）}

\begin{itemize}
  \item GitHub: \url{https://github.com/vulca-org/vulca-visual-control-sdk}
  \item Vulca Plugin: \url{https://github.com/vulca-org/vulca-visual-agent-plugin}
  \item PyPI: \url{https://pypi.org/project/vulca/}
  \item 截至 2026-04-20 README refresh 节点：累计 \keyword{82 views / 31 unique visitors / 6 stars}（基线），
  正在进入种子内容驱动的流量窗口。后续数据在每季度复盘中更新。
  \item 社区种子内容：dev.to（英文 deep-dive）、Xiaohongshu 中文 carousel（Bieber / Trump 案例 2026-04-18 发布）、
  GitHub Discussions 首批讨论。
\end{itemize}

\subsection*{D. 学术产出与合作}

\begin{itemize}
  \item \keyword{ACM MM 2026 EmoArt Challenge（参赛）}\quad 团队以 Vulca 栈作为 Track 1 \textit{Emotion-Aware Artistic Image Generation} 参赛方案——\keyword{为项目提供学术与工业双边的第三方背书}；Track 2 涉及情绪识别子方向的后续研究输出；提交截止 2026-06-10。
  \item \keyword{\code{docs/superpowers/} 技术文档体系}\quad 所有设计、brainstorming、spec、plan、ship-gate 公开在仓库，
  构成自证的方法论资产。
\end{itemize}

\subsection*{E. Provider 矩阵详细说明}

\begin{center}
\small
\begin{tabularx}{\linewidth}{@{}>{\tabRR}p{2.2cm}>{\tabRR}p{3.0cm}>{\tabRR}X@{}}
\toprule
\rowcolor{tableHead}
\heiti Provider & \heiti 默认模型 & \heiti 适用场景 \\
\midrule
\code{comfyui} & SDXL / FLUX / 自定义 workflow & 本地离线；自主可控首选 \\
\code{gemini} & \code{gemini-3.1-flash-image-preview} & 云端高质量；credits 覆盖 \\
\code{gemini} (VLM) & \code{gemini-2.5-pro} via LiteLLM & 视觉理解、评估推理辅助 \\
\code{openai} & DALL·E 系列 & 云端兼容；特定客户偏好 \\
\code{mock} & 无模型调用 & CI 与本地开发 \\
\bottomrule
\end{tabularx}
\end{center}

\subsection*{F. 术语表}

\begin{description}[leftmargin=5.6em,style=nextline,itemsep=2pt]
  \item[\keyword{Agent-native}] 工具的接口与返回值为 agent 规划设计，不是为人类 UI 设计。
  \item[\keyword{MCP}] Model Context Protocol，由 Anthropic 于 2024 年提出，是 agent 与外部工具交互的开放协议。
  \item[\keyword{Skill}] Claude Code 中的一级扩展单位，包含声明式 prompt + 工具调用脚本；可被 agent 自动调用。
  \item[\keyword{Ship-gate}] Vulca 的 pre-release 验证闸门：simulated + live MCP 双层验证，全绿才允许发布。
  \item[\keyword{L1–L5}] Vulca 文化评估体系的五层评价指标（形制 / 技法 / 语汇 / 风格 / 气韵）。
  \item[\keyword{Provider}] 底层模型访问抽象层，支持 comfyui / gemini / openai / mock 等；可 runtime 切换。
  \item[\keyword{EVF-SAM}] 基于 SAM 的语言引导分割模型，Vulca 语义分层核心组件之一；MPS 下约 3s / prompt。
\end{description}

\subsection*{G. 关键问题与答复}

以下问答覆盖投资人与申报评审关注的核心疑问，同步梳理 Vulca 当前阶段的立场与路线。

\begin{enumerate}
  \item \keyword{Q：从目前 6 stars / 82 views 到 Q4 2026 1500 stars + 30 付费账户，这中间是怎么发生的？}\\
  A：漏斗见第 \ref{sec:roadmap} 章（7.4 节）——曝光 / 知晓 / 尝试 / 激活 / 留存 / 付费意向 / 付费转化七级；配合 dev.to 月度深度文 + Xiaohongshu 双语种子内容 + Claude Code 插件生态生态位 + ACM MM 2026 学术引流。如果某一级转化不达预期，优先加码的渠道组合：\keyword{dev.to 付费推广 + 针对 agent 创业项目的 Discord/Slack 置换投放 + 1-2 场线上 AI 开发者 meetup 主题分享}（整体以内容驱动为主、付费推广为辅），必要时降级目标。
  \item \keyword{Q：文化评估 L1–L5 YAML 能否被 ComfyUI 社区 / 头部厂商两个月内复刻？}\\
  A：理论上可以复制\keyword{格式}，但难以复制(1) 与学术合作方一起迭代的\keyword{规则原创性}、(2) 跨 13 种传统的\keyword{覆盖宽度}、(3) 社区扩展的\keyword{数据资产飞轮}、(4) 已经进入 ACM MM 2026 的\keyword{学术背书链}。真实护城河 = 时间 + 合作关系网络 + 贡献者生态，不是代码本身。
  \item \keyword{Q：Google Gemini / Anthropic 的 credits 到期后怎么办？}\\
  A：参见第 \ref{sec:product}.5 节和第 \ref{sec:financials}.2 节——credits 只覆盖 demo 与有限早期使用；生产路径由\keyword{本地 ComfyUI + Ollama}\textbf{兜底}，一键切换；客户侧由 pricing 设计传导部分成本；需要时可与上游厂商续期议价。不依赖任何单一厂商。
  \item \keyword{Q：首批目标客户画像与接触策略是什么？}\\
  A：目标客户画像包括博物馆数字化、出版教辅、品牌设计集团、高校艺术研究中心四类；本轮募资交割后进入正式 BD 阶段，优先通过学术合作方（Dundee DJCAD / ACM MM 学术圈）和 Xiaohongshu / dev.to 的自然流量引入客户线索。具体接触清单随 BD 负责人到岗后公开于数据室。
  \item \keyword{Q：团队目前 solo，本轮融资后如何保证 ship 频率不下滑？}\\
  A：参见第 \ref{sec:team} 章和第 \ref{sec:risks} 章人才部分——(1) 工程流程已标准化（brainstorm → spec → plan → ship gate，整套开源于 \code{docs/superpowers/}，任何新成员可按图索骥）；(2) 天使轮第一笔用途就是 de-risk 关键人依赖（P0 岗位为资深算法 + 产品/BD）；(3) 每个季度有明确的 shipped skill 指标为硬约束，新成员完成度可量化。
\end{enumerate}

\subsection*{H. 创始人简历（Founder CV）}

\begin{center}
\begin{tabularx}{\linewidth}{@{}>{\tabRR}p{2.8cm}>{\tabRR}X@{}}
\toprule
\rowcolor{tableHead}
\multicolumn{2}{@{}l}{\heiti\color{vulcaPrimary}\large Yu Haorui} \\
\midrule
联系方式 & yuhaorui48@gmail.com \textperiodcentered{} 2655435@dundee.ac.uk \\
GitHub & \url{https://github.com/yha9806} \textperiodcentered{} org \url{https://github.com/vulca-org} \\
角色 & Vulca 创始人 \& 技术负责人 \\
\midrule
\multicolumn{2}{@{}l}{\heiti\color{vulcaPrimary}教育背景} \\
\midrule
2024 – 至今 & \textbf{PhD Researcher}, University of Dundee, Duncan of Jordanstone College of Art \& Design (DJCAD) \newline
导师：Prof.~Natasha Lushetich（DJCAD）\newline
合作者：Dr.~Ramon Ruiz-Dolz（Computing 系 Lecturer, ELLIS Society）\newline
论文题目：\textit{Evaluating Vision-Language Models for Cultural Understanding in Heritage and Arts Contexts} \\
\midrule
\multicolumn{2}{@{}l}{\heiti\color{vulcaPrimary}代表论文} \\
\midrule
EMNLP 2025 Findings & \textbf{一作} \textperiodcentered{} \textit{A Structured Framework for Evaluating and Enhancing Interpretive Capabilities of Multimodal LLMs in Culturally Situated Tasks} \textperiodcentered{} 与 R. Ruiz-Dolz、Q. Yi 合作 \textperiodcentered{} pp.~1945–1971 \textperiodcentered{} \url{https://aclanthology.org/2025.findings-emnlp.103/} \\
ACL 2026（Dataset Track 审稿中） & \textbf{一作} \textperiodcentered{} \textit{VULCA-Bench: A Benchmark for Culturally-Aware Visual Understanding at Five Levels} \textperiodcentered{} arXiv 阶段性快照 8 种 / 7410 样本；SDK 当前已扩展至 13 种 \textperiodcentered{} \href{https://arxiv.org/abs/2601.07986}{arXiv 2601.07986} \\
arXiv 2026 & \textbf{一作} \textperiodcentered{} \textit{Cross-Cultural Expert-Level Art Critique Evaluation with Vision-Language Models} \textperiodcentered{} 与 X. Wen、F. Zhang、Q. Yi 合作 \textperiodcentered{} \href{https://arxiv.org/abs/2601.07984}{arXiv 2601.07984} \\
\midrule
\multicolumn{2}{@{}l}{\heiti\color{vulcaPrimary}开源 / 工程产出} \\
\midrule
Vulca SDK & 独立设计与实现 \textperiodcentered{} Apache 2.0 \textperiodcentered{} v0.14.0 (2026-04-09) $\to$ v0.17.5 (2026-04-23)，累计 11 个正式版本，平均每 1.3 天一次 ship \\
技术基线 & 21 个 MCP 工具 \textperiodcentered{} 3 个端到端 shipped skill（\code{/decompose}、\code{/visual-brainstorm}、\code{/visual-spec}）\textperiodcentered{} 高覆盖率测试套件，CI 绿灯率 $\ge$ 99\% \\
工程方法论 & brainstorm $\to$ spec $\to$ plan $\to$ ship-gate 完整流程，开源于 \code{docs/superpowers/} 目录 \\
生态位 & Claude Code 插件市场 \code{vulca-org/vulca-visual-agent-plugin}；PyPI \code{vulca} \\
\midrule
\multicolumn{2}{@{}l}{\heiti\color{vulcaPrimary}学术与资源网络} \\
\midrule
主要学术顾问与研究资源支持方 & Prof.~Jindong Wang（William \& Mary 助理教授，前 MSRA 资深研究员 2019–2024，23000+ citations, H-index 54，World Top 2\% Highly Cited Scientists）—— 已建立长期学术合作关系，书面材料备于 DD 数据室 \\
共同作者 & Dr.~Ramon Ruiz-Dolz（Dundee Computing），Qiufeng Yi（Birmingham Mechanical Engineering） \\
模型厂商关系 & 已获 Google Gemini + Anthropic 研究级 / 创业级 API credits \\
学术赛事 & ACM MM 2026 EmoArt Challenge Track 1 \textit{Emotion-Aware Artistic Image Generation} 参赛（以 Vulca 栈为技术方案；提交 2026-06-10；为项目提供第三方学术背书） \\
\midrule
\multicolumn{2}{@{}l}{\heiti\color{vulcaPrimary}研究方向与兴趣} \\
\midrule
方向 & 视觉-语言模型的文化评估 \textperiodcentered{} agent-native 多模态工作流 \textperiodcentered{} 开源生态与工程方法论 \\
个人驱动 & 把"AI 理解中国艺术"从 benchmark 变成可交付产品，让文化语境感知成为每个 AI agent 的默认能力 \\
\bottomrule
\end{tabularx}
\end{center}


